% =====================================================================
% CAPITULO 6 - DISCUSION
% Sustituye al capitulo 6 de la version anterior de la memoria.
% =====================================================================

\section{Discusión}
\label{cap:discusion}

\subsection{Reorientación del objeto de estudio}
\label{sec:disc-reorientacion}

El planteamiento inicial de este trabajo perseguía identificar biomarcadores
transcriptómicos universales de cáncer de pulmón mediante la integración de
múltiples cohortes públicas. Los resultados del capítulo anterior no sostienen
ese objetivo, y conviene ser explícito sobre por qué.

La firma derivada del contraste entre tejido tumoral y tejido sano adyacente
reúne tres características que impiden interpretarla como un conjunto de
biomarcadores: su señal procede en parte sustancial de la composición del tejido
más que de biología tumoral específica (sección~\ref{sec:composicion}); su
estabilidad aparente era un artefacto del diseño de validación
(sección~\ref{sec:falacia}); y la tarea que resuelve no constituye un problema
clínico, dado que la distinción entre tumor y parénquima normal la resuelve el
examen histopatológico convencional con fiabilidad superior, coste
incomparablemente menor y sobre la misma muestra ya disponible.

Esto no invalida el framework desarrollado. Lo que invalida es la lectura que se
había hecho de sus resultados. El pipeline descarga, normaliza, curó
metadatos mediante modelos de lenguaje, ejecuta análisis diferencial, entrena
modelos y valida de forma externa; todos esos componentes funcionan, y el
control positivo de la sección~\ref{sec:subtipo} demuestra que recuperan biología
reproducible cuando está presente. La contribución defendible de este trabajo no
es, por tanto, una firma molecular, sino la caracterización sistemática de los
modos en que estos pipelines fallan sin avisar, y un conjunto de controles
concretos para detectarlos.

\subsection{Los cuatro modos de fallo silencioso}
\label{sec:disc-fallos}

Los problemas documentados en la sección~\ref{sec:auditoria} comparten un rasgo
que los hace especialmente peligrosos: ninguno interrumpe la ejecución. El
pipeline termina, escribe sus ficheros y produce tablas de aspecto correcto.

\paragraph{El desalineamiento como fallo indistinguible.} El caso de GSE30219 es
instructivo porque su efecto imita exactamente el de una hipótesis nula
verdadera. Un AUC de 0,56 sobre una cohorte de 146 tumores es perfectamente
compatible con la conclusión «esta señal no generaliza a esta cohorte», y ese es
el tipo de resultado que un trabajo honesto acepta y reporta. Solo la
comparación con marcadores de referencia externos --- la observación de que KRT5
mostraba una diferencia de $+0{,}90$ en esta cohorte frente a $+5{,}32$ y
$+5{,}66$ en las otras dos --- reveló que la señal estaba presente pero
desacoplada de las etiquetas. La lección metodológica es que la validación
frente a conocimiento biológico previo no es un adorno interpretativo, sino un
control de integridad de los datos.

Resulta pertinente señalar que el análisis original no se vio afectado por este
error, y la razón es instructiva: GSE30219 contiene una sola clase, de modo que
intercambiar etiquetas idénticas entre sí no altera nada. El fallo permaneció
latente y solo se manifestó al plantear una pregunta que utilizara la
información clínica. Un pipeline puede arrastrar errores de este tipo durante
todo su ciclo de vida sin que se manifiesten, hasta que alguien formula la
pregunta adecuada.

\paragraph{La curación automatizada como fuente de sesgo de selección.} El
empleo de modelos de lenguaje para homogeneizar metadatos clínicos heterogéneos
es una de las aportaciones metodológicamente más interesantes del framework, y
su comportamiento bimodal --- éxito prácticamente completo en ocho cohortes,
colapso en dos --- merece una lectura cuidadosa. El problema no es la tasa
global del 82,8\,\%, sino que las pérdidas no son aleatorias: se concentran en
cohortes concretas y, dentro de ellas, afectan a la práctica totalidad de las
muestras. GSE40791 pasó de 194 muestras a 12, y las 12 supervivientes presentan
una composición de clases (10 controles y 2 tumores) que no representa al
estudio original. El resultado es un sesgo de selección introducido por una capa
de procesamiento automático, no declarado como tal en ningún punto del
pipeline. La recomendación que se deriva es que cualquier etapa de curación
automatizada debe reportar su tasa de éxito por cohorte como parte de los
resultados, y que las cohortes con tasas bajas deben excluirse explícitamente en
lugar de participar con su fracción superviviente.

\paragraph{Cohortes de una sola clase.} Que GSE30219 y GSE50081 no contengan
controles no es un defecto de esas cohortes: son estudios de supervivencia,
diseñados para correlacionar expresión con evolución clínica, y su composición es
la correcta para su propósito. El error consiste en emplearlas como conjunto de
test de un clasificador binario. Este fallo es además el más fácil de detectar
automáticamente --- basta comprobar el número de clases presentes antes de
evaluar --- y el más fácil de pasar por alto cuando la métrica reportada es
únicamente la \emph{accuracy}, porque produce precisamente los valores más altos
de la tabla.

\subsection{Discriminación y decisión como propiedades independientes}
\label{sec:disc-calibracion}

El hallazgo de la sección~\ref{sec:calibracion} no estaba previsto y es, a
nuestro juicio, el resultado de mayor utilidad práctica del trabajo. La
coexistencia de un AUC medio de 0,925 con una \emph{balanced accuracy} media de
0,773, y su expresión extrema en GSE40791 (\emph{accuracy} 0,167 con AUC 1,000),
descompone lo que habitualmente se trata como un único atributo --- «el modelo
generaliza» o «no generaliza» --- en dos propiedades que transfieren de forma
independiente entre cohortes.

La capacidad de ordenación transfiere bien: la firma sitúa consistentemente las
muestras tumorales por encima de los controles en cohortes que nunca ha visto. El
umbral de decisión no transfiere: el punto de corte aprendido sobre una
distribución de clases determinada (938 tumores frente a 219 controles) deja de
ser adecuado cuando la cohorte de destino presenta otra composición. La
sensibilidad media de 0,983 frente a una especificidad media de 0,563 describe un
modelo que ha aprendido, razonablemente dada su dieta de entrenamiento, que
predecir «tumor» casi siempre acierta.

Esta lectura reinterpreta un patrón que el análisis original registraba sin
explicar. Los casos de \emph{accuracy} próxima a 0,5 acompañada de AUC próximo a
1,0 se habían atribuido a la dificultad intrínseca del problema, al tamaño
muestral reducido o a particularidades técnicas de los estudios. No era ninguna
de esas cosas: era un problema de calibración, y los problemas de calibración
son corregibles. La vía inmediata es la ponderación de clases en el
entrenamiento; la más robusta, reportar métricas independientes del umbral (AUC,
precisión media) junto con métricas de decisión, y ajustar el punto de corte
sobre la cohorte de destino cuando el escenario de aplicación lo permita.

\subsection{Qué mide la firma: composición, proliferación y biología}
\label{sec:disc-composicion}

La hipótesis de que la firma tumor frente a sano capturase esencialmente
composición tisular no se confirmó al umbral fijado de antemano
($|\rho|>0{,}7$; obtenido $\rho=-0{,}637$). Mantener ese criterio obliga a una
conclusión más matizada, y probablemente más correcta, que la hipótesis de
partida.

La composición tisular explica una fracción sustancial de la señal. La evidencia
es convergente: la correlación entre el \emph{score} del clasificador y el
contenido de pulmón normal persiste al restringir el análisis a muestras
tumorales, alcanzando $\rho=-0{,}870$ ($p=7\times10^{-71}$) en la cohorte con
mayor número de tumores evaluables; y cinco de los cincuenta genes principales de
la firma son marcadores canónicos de epitelio alveolar y estroma pulmonar normal,
con $\log$FC en torno a $-4$. Interpretados como biomarcadores tumorales, genes
como AGER, CLDN18 o SFTPC informan principalmente de cuánto parénquima sano
conserva la muestra.

Pero no la agotan. El segundo eje detectado --- correlación positiva con
proliferación, entre $+0{,}267$ y $+0{,}719$ --- indica que el clasificador
integra al menos dos fenómenos distintos, y el segundo sí es biología tumoral en
sentido estricto. La descripción más ajustada de lo que mide la firma es
«pérdida de arquitectura alvéolo-capilar normal más ganancia de actividad
proliferativa», una combinación real, reproducible y de escaso valor
discriminativo por ser común a la práctica totalidad de los tumores sólidos.

Debe señalarse una limitación de este análisis: solo cuatro cohortes reunían
tumores suficientes ($n\geq8$) para el contraste restringido, y los estimadores
presentan heterogeneidad notable entre ellas ($\rho$ de $-0{,}343$ a
$-0{,}870$). Un abordaje más riguroso emplearía métodos de deconvolución celular
para estimar la proporción de cada tipo celular, en lugar de un panel de
marcadores promediado.

\subsection{Implicaciones sobre el diseño de la validación externa}
\label{sec:disc-lodo}

La estrategia LODO es metodológicamente superior a la validación cruzada
convencional dentro de un mismo estudio, y su adopción sigue siendo escasa en la
literatura de biomarcadores transcriptómicos. Este trabajo no cuestiona esa
elección, sino dos usos concretos que se hicieron de ella.

El primero es de terminología con consecuencias conceptuales: LODO no corrige el
efecto lote, lo \emph{mide}. Presentarla como mecanismo de superación del
\emph{batch effect} confunde la naturaleza del procedimiento. En el pipeline
analizado no existe corrección de lote alguna --- solo normalización por
cuantiles dentro de cada estudio y estandarización por estudio ---, y la
degradación de rendimiento entre cohortes que LODO revela \emph{es} la magnitud
del problema, no su solución.

El segundo es el uso de los \emph{folds} LODO como réplicas para medir
estabilidad. La demostración de la sección~\ref{sec:falacia} es cuantitativa y,
en nuestra opinión, generalizable más allá de este trabajo: cuando el número de
cohortes es moderado, los conjuntos de entrenamiento de los distintos
\emph{folds} comparten una mediana del 97,6\,\% de sus muestras, y la
concordancia direccional resultante es del 100,0\,\% en la totalidad de los
pares. Una métrica cuyo valor está determinado por el diseño no puede aportar
evidencia sobre el fenómeno estudiado. El control con mitades disjuntas de
tamaño comparable (73,98\,\%) delimita cuánto de ese 100\,\% procede del
solapamiento.

La consecuencia sobre los siete genes previamente destacados es contundente:
ninguno replica al ajustar modelos independientes. Esto no significa que carezcan
de interés biológico --- SLC6A4 y KANK3 son marcadores endoteliales pulmonares
bien caracterizados, coherentes con la interpretación composicional de la
sección anterior ---, sino que la evidencia aportada para priorizarlos no era
evidencia. La alternativa empleada en la sección~\ref{sec:subtipo}, exigir
concordancia de tamaño de efecto en cohortes ajustadas por separado, es más
exigente y produjo una firma que recupera marcadores de uso clínico establecido.

\subsection{Alcance y limitaciones del control positivo}
\label{sec:disc-subtipo}

La clasificación de subtipo histológico cumple su función como control positivo:
demuestra que el marco metodológico detecta biología reproducible cuando existe,
con AUC medio de 0,967 y recuperación completa de los doce marcadores de
inmunohistoquímica diagnóstica. Conviene, no obstante, precisar tres límites.

Primero, la distinción entre adenocarcinoma y carcinoma escamoso está
extensamente caracterizada en transcriptómica y su recuperación no constituye una
aportación original. Su valor es de verificación, no de descubrimiento.

Segundo, el resultado se obtuvo excluyendo el 31,3\,\% de los tumores
disponibles, precisamente los de histología ambigua. El análisis de la
sección~\ref{sec:dificiles} muestra que el modelo es más prudente ante esas
muestras de lo previsto (29,9\,\% de asignaciones de alta confianza frente al
62,6\,\% en muestras de clase conocida; la hipótesis registrada no se confirmó),
pero también que el 93\,\% de los 101 tumores neuroendocrinos se etiqueta como
adenocarcinoma. Dado que el carcinoma microcítico requiere un abordaje
terapéutico completamente distinto, esta es una limitación con consecuencia
clínica y no meramente estadística.

Tercero, las tres cohortes comparten plataforma (GPL570), de modo que el
resultado no informa sobre transferencia entre tecnologías de medición. Las tres
cohortes de RNA-seq que habrían permitido evaluarlo son precisamente las que el
pipeline no llegó a procesar.

\subsection{Recomendaciones metodológicas}
\label{sec:disc-recomendaciones}

De los resultados obtenidos se derivan seis controles concretos, aplicables a
cualquier pipeline de descubrimiento de biomarcadores sobre repositorios
públicos:

\begin{enumerate}
\item \textbf{Alinear muestras y metadatos por identificador explícito}, nunca
      por posición, y verificar la correspondencia como paso previo a cualquier
      análisis. La coincidencia de longitudes no implica coincidencia de orden.
\item \textbf{Validar contra marcadores biológicos conocidos} antes de
      interpretar cualquier resultado. Un panel de referencia de la patología
      estudiada funciona como control de integridad de los datos y detecta
      fallos que ninguna comprobación de formato revela.
\item \textbf{Reportar el \emph{baseline} de clase mayoritaria} junto a toda
      métrica de rendimiento, y excluir explícitamente las cohortes de una sola
      clase de la evaluación, indicándolo en las tablas.
\item \textbf{Separar métricas de discriminación y de decisión.} Informar AUC y
      \emph{balanced accuracy} conjuntamente, con sensibilidad y especificidad
      desglosadas; su divergencia identifica problemas de calibración que la
      \emph{accuracy} agregada oculta.
\item \textbf{Medir la estabilidad sobre particiones disjuntas}, no sobre
      \emph{folds} solapados, y declarar la fracción de muestras compartidas
      cuando se empleen estas últimas.
\item \textbf{Cuantificar y reportar la tasa de éxito de toda etapa de curación
      automatizada} por cohorte, tratando las pérdidas elevadas como criterio de
      exclusión y no como merma silenciosa del tamaño muestral.
\end{enumerate}

\subsection{Líneas futuras}
\label{sec:disc-futuro}

Los datos ya descargados permiten abordar preguntas que este trabajo no ha
explorado y que resultan más informativas que la comparación entre tumor y tejido
sano. GSE30219 y GSE50081 disponen de seguimiento clínico completo (307 y 181
tumores, con 200 y 75 eventos de mortalidad respectivamente), lo que habilita
análisis de supervivencia con validación cruzada entre cohortes. GSE31210
incluye el estado de los principales genes conductores --- 127 tumores con
mutación de EGFR, 20 de KRAS, 11 con fusión de ALK y 68 sin ninguna de las tres
alteraciones ---, lo que permite plantear la predicción de alteración
conductora a partir del perfil de expresión. Ambas son preguntas cuya respuesta
no se conoce de antemano por inspección de la muestra, condición que la tarea
tumor frente a sano no cumplía.

En el plano metodológico, tres extensiones mejorarían el rigor de los análisis
presentados: el empleo de deconvolución celular en lugar de paneles de
marcadores promediados para cuantificar composición tisular; la incorporación de
corrección de lote explícita, con evaluación de su efecto sobre la transferencia
entre cohortes; y la recuperación de las tres cohortes de RNA-seq no procesadas,
que permitirían por primera vez evaluar la generalización entre plataformas de
medición distintas.
